\chapter[Additional Reproducibility Notes]{Additional Test Logs and Reproducibility Notes}
\label{app:reproducibility-notes}

\section{Clean measurement discipline}

The gas measurements in this report were interpreted under the following rules:

\begin{itemize}
  \item never compare a fixed infrastructure cost with a marginal per-exchange cost;
  \item never compare a local Step 8 cost with a complete dispute total without saying so;
  \item restart or reset the Hardhat chain when deployment addresses become stale;
  \item rebuild native CLIs before running native large-file benchmarks;
  \item report whether a number comes from the initial, monolithic, or final specialized
  implementation.
\end{itemize}

These rules were introduced after early measurements mixed scopes and were difficult to
explain. They are now part of the experimental methodology.

\section{Canonical benchmark commands}

\begin{verbatim}
cd src/hardhat

npx hardhat test \
  test/performance/HanaComparisonCurrent.test.ts
npx hardhat test \
  test/performance/SpecializedContractArchitecture.test.ts
npx hardhat test \
  test/performance/SpecializedFinalSameScope.test.ts
\end{verbatim}

For the native hardcoded benchmark:

\begin{verbatim}
cd src/wasm
cargo build --release --bin precontract_cli \
  --bin hardcoded_sha256_dispute_cli

cd ../hardhat
npx hardhat test \
  test/performance/NativeHardcodedSHA256FullDispute1GB.test.ts
\end{verbatim}

For a small smoke test of the same code path:

\begin{verbatim}
SOX_FULL_DISPUTE_SIZE_BYTES=16384 npx hardhat test \
  test/performance/NativeHardcodedSHA256FullDispute1GB.test.ts
\end{verbatim}

\section{Interpreting local benchmark output}

Hardhat gas values come from transaction receipts and are stable for a fixed contract
version, optimizer configuration, and test path. Wall-clock times are less stable because
they depend on the host machine, file system, CPU load, and whether the benchmark is run
inside WSL or directly on Linux.

The 1 GiB benchmark creates a temporary sparse file through truncation. This makes the test
practical without storing a permanent 1 GiB fixture in the repository. The benchmark still
executes the precontract CLI and dispute CLI on that file and then performs the on-chain
dispute path.

\section{Known local-environment issues}

\begin{itemize}
  \item If \texttt{run-alto.sh} reports that the EntryPoint does not exist, the Hardhat
  chain was probably restarted after deployment. Redeploy with \texttt{./deploy-all.sh}.
  \item If a native CLI is missing, rebuild it from \path{src/wasm} with
  \texttt{cargo build --release --bins}.
  \item If Hardhat cannot write under \path{src/hardhat/artifacts}, the files may have been
  created by another user. Restore ownership before recompiling.
  \item If the web application has stale exchanges, reset \path{src/app/db/sox.sqlite} and
  replay \path{src/app/db/init.sql}.
\end{itemize}

\section{Submitted-code boundary}

The repository contains historical code, intermediate monolithic contracts, final
specialized contracts, scripts, and benchmark files. The report uses the following
terminology consistently:

\begin{itemize}
  \item initial implementation: the baseline inherited from the previous project;
  \item intermediate monolithic implementation: the variant-rich contracts before the final split;
  \item final specialized implementation: factory, clones, specialized disputes, and the
  hardcoded strict path.
\end{itemize}

This boundary is important for grading and future reuse. The project contribution is not
only the final gas number; it is also the separation of these stages, the implementation of
the retained variants, and the reproducible measurement framework.
